%%%%%%%%%%%%%%%%%%%%%%%%%%%%%%%%%%%%%%%%%%%%%%%%%%%%%%%%%%%%%%%%%%%%%%%%%
%
%	LaTeX File for Stanford University PhD Thesis
%
%%%%%%%%%%%%%%%%%%%%%%%%%%%%%%%%%%%%%%%%%%%%%%%%%%%%%%%%%%%%%%%%%%%%%%%%%
%	Copyright 2001  by  Jung-Suk Goo    (goojs@gloworm.stanford.edu)
%%%%%%%%%%%%%%%%%%%%%%%%%%%%%%%%%%%%%%%%%%%%%%%%%%%%%%%%%%%%%%%%%%%%%%%%%

\beforepreface
\prefacesection{Abstract}
{\noindent\fancypar{T}{{\sc he}} development of foundation models for wireless communications requires two key components: a powerful signal encoder to extract rich features and an intelligent reasoning model for coherent, cross-task decision-making. As building such models from scratch is prohibitively expensive, this dissertation presents a practical framework that creates robust wireless foundation models by strategically approximating and integrating world knowledge.\\
Our cost-effective approach proposes a modular, three-component framework and provides proof-of-concept validation for each component. The first is a rich feature encoder, for which we develop DenoMAE. Through unsupervised learning on unlabeled signals, these models approximate world knowledge of universal signal representations by treating noise as an informative modality and processing multimodal time-domain and constellation-diagram data. The second is an intelligent reasoning engine, for which we demonstrate a novel proof-of-concept by utilizing LLMs to perform open-set classification based on statistical signal features, successfully transferring their pre-trained world knowledge to the wireless domain. The third is an adaptation module, for which we employ Meta-Task learning to enable the framework to rapidly apply its integrated world knowledge to new tasks and channel conditions with minimal labeled data.\\
This dissertation demonstrates that the key to economical and effective wireless foundation models lies not in massive, isolated training, but in the design of this strategic, modular framework. By validating the components for unsupervised representation (DenoMAE), LLM-based reasoning, and meta-learning for adaptation, our work provides a practical pathway and a conceptual blueprint for achieving foundation model capabilities in wireless communications.}

% \fancypar{W}{{\sc e }} propose novel approaches to representation learning and robust classification that enable generalization in data-constrained environments, specifically addressing challenges posed by limited labeled data, noisy inputs, and distribution shifts. Through a series of interconnected studies, we demonstrate that carefully designed representation learning frameworks can significantly improve performance across diverse domains including wireless communication, image classification, and few-shot learning. Our contributions span three primary research directions: (1) noise-robust representation learning for image classification and wireless communication signals, (2) episodic representation learning for meta-learning, and (3) efficient utilization of unlabeled data through self-supervised, semi-supervised and few-shot learning approaches. The proposed methodologies consistently demonstrate superior performance over existing approaches, particularly in low-data regimes and under challenging noise conditions, offering practical solutions for real-world applications where data acquisition is costly or limited.  